{\large\textbf{Maxwell, Rayleigh and Mount Everest: THE PROBLEM}}\footnote{Amitabh Virmani (CMI, Chennai) and A. C. Biyani (retired Govt. Nagarjuna P.G. College of Science. Raipur) were the principal authors of this problem. The contributions of the Academic Committee, Academic Development Group, and the International Board are gratefully acknowledged.}

Lord Rayleigh visited Darjeeling, India in 1897. On viewing Mount Everest at a
distance of 170 km he was reminded of a query by James Maxwell 26 years
earlier on the attenuation of light by air and the visibility of far-off
peaks. Lord Rayleigh authored a famous research
paper\footnote{Philosphical Mag. ``On the transmission of light through an atmosphere ... and the origin of the blue of the sky'', Vol. 47, pg 375-384 1899}
two years later in 1899 in which he investigated this problem. In the present
problem we attempt to reconstruct, at least partially, a modern version of
Lord Rayleigh's line of reasoning.

\textbf{Oscillation of the electron cloud:} We model a typical neutral air
molecule by a stationary positive charge $q$ surrounded by a spherical uniform
charge cloud of mass $m$, radius $r$ and charge $-q$. The natural angular
frequency of vibration of the molecule is $\omega_0$. Light is incident on it
and the negative cloud oscillates maintaining its spherical shape with angular
frequency $\omega$ as,
\begin{equation}
y = y_0\cos(\omega t), \tag{1}
\end{equation}
under the influence of the electric field
\begin{equation}
\vec{E}(t) = E_0\cos(\omega t)\,\hat{y}, \tag{2}
\end{equation}
of the light wave. Here $y$ represents the separation between the stationary
positive charge and the centre of the negative charge cloud of the molecule.

\textbf{A.1}\quad Set up the equation for the motion for $y$. Take the
contribution to the magnitude of the acceleration due to the electric field to
be $E(t)\,q/m$. \hfill 0.5pt

\textbf{A.2}\quad Based on the information provided above, solve the equation
for $y$. Find the amplitude. \hfill 0.5pt

\textbf{A.3}\quad Find the magnitude $p(t)$ of the dipole moment of the air
molecule as a function of time for $\omega \ll \omega_0$. \hfill 0.5pt

\textbf{A.4}\quad Obtain an expression for $\omega_0$ in terms of $q$, $m$ and
$r$. \hfill 0.5pt

\textbf{Power radiated:} The sinusoidal time dependent dipole radiates
electromagnetic radiation. The power radiated depends on the amplitude of the
dipole moment $p_0 = qy_0$, the permittivity of vacuum $\epsilon_0$, the speed
of light $c$, and the frequency of oscillation $\omega$.

\textbf{B.1}\quad Use dimensional analysis to express the average power $s$
radiated in terms of these quantities. \hfill 1pt

\textbf{B.2}\quad Take the proportionality constant to be $1/12\pi$. Express
the answer for the power radiated $s$ in terms of $E_0$, $\omega_0$, $\omega$
and related quantities. Take $\omega \ll \omega_0$. \hfill 0.2pt

\textbf{Attenuation of the Intensity} $I(x)$\,\textbf{:} Recall that the
intensity of EM waves is
\begin{equation}
\frac{1}{2}c\epsilon_0E_0^2. \tag{3}
\end{equation}
The intensity decreases along the path of light because the power
\begin{equation}
S = n_0s \tag{4}
\end{equation}
is lost per unit volume. Here $n_0$ is the number of molecules per unit
volume.

\textbf{C.1}\quad Set up the differential equation for the intensity $I(x)$ as
a function of the distance $x$. \hfill 1pt

\textbf{C.2}\quad Obtain the expression for the intensity $I(x)$ as a function
of $x$ in terms of the characteristic length scale $L$ associated with the
drop in intensity. Take the initial intensity to be $I_0$. \hfill 0.5pt

\textbf{C.3}\quad Take $m$ to be the mass of the electron (typically only one
electron constitutes the charge cloud) and take
\begin{align}
&n_0 = 2.54\times10^{25}\ \mathrm{m}^{-3}, \tag{5}\\
&\omega_0 = 1.25\times10^{16}\ \mathrm{rad}\cdot\mathrm{s}^{-1}, \tag{6}\\
&\omega = 3.25\times10^{15}\ \mathrm{rad}\cdot\mathrm{s}^{-1}. \tag{7}
\end{align}
Obtain the numerical value of $L$ in kilometers. \hfill 0.3pt

\textbf{D.1}\quad \textbf{Height} $H'$ \textbf{of the Mountains as seen by an
observer\,:} In the figure the point $P$ denotes Darjeeling, a hill station in
the eastern Himalayas at height $h$ = 2042 m above the sea level. The line
$BS$ denotes Mt Everest which is $d$ = 170 km away from Darjeeling and is of
height $H$ = 8848 m. Another peak Mount Kanchenjunga (not shown in the figure)
is 75 km away from Darjeeling and is of height 8586 m. Obtain an expression
and the numerical values of the vertical height $H'$ of these mountains as
seen by an observer from Darjeeling in terms of the above-mentioned
quantities. Assume that the observer is unable to see below the local
horizon. Draw an appropriate figure. Take the radius $R$ of the Earth to be
6378 km. \hfill 2pt

\begin{center}\includegraphics[width=0.5\linewidth]{figures/APhO_2022_Q3_fig1.png}\end{center}

\begin{center}Figure 1. Great circle on which lie the mountain $BS$ at height
$H$ and the observer $P$ at height $h$. Note that the figure is not to
scale.\end{center}

\textbf{E.1}\quad Take the intensity of Mt Kanchenjunga at Darjeeling to be
the reference value. What would be the intensity of Mt Everest relative to Mt
Kanchenjunga as seen from Darjeeling? In this problem we ignore the variation
in the number density of air molecules with height. If the intensity is 5 \%
or more of the reference value the mountain is said to be visible. Will Mt
Everest be visible from Darjeeling? \hfill 1pt

\textbf{F.1}\quad \textbf{Attenuation length} $L_p$ \textbf{due to aerosol
pollution\,:} Above, we calculated the characteristic length scale $L$
associated with the drop in intensity due to scattering with air molecules. We
are now interested in the characteristic length $L_p$ associated with the drop
in intensity due to scattering with \textit{aerosol} particles (pollution).
$L_p$ depends on the number density $n$ of particles and the cross sectional
area $\pi r^2$ of the aerosol particle of radius $r$. Obtain this dependence
using physical insight and dimensional analysis. Take the dimensionless
constant to be 1/8. Given mild pollution the average aerosol density at
Darjeeling is $\rho_p = 5\,\mu\mathrm{g/m^3}$, and their average radius is 500
nanometres. What is the length $L_p$? Let the density of an individual aerosol
particle be $\rho = 3\,\mathrm{g/cm^3}$. Note 1 $\mu$g = 10$^{-9}$ kg and 1 nm
= 10$^{-9}$ m. \hfill 1pt

\textbf{G.1}\quad \textbf{Relative intensity and Visibility of Mt.
Kanchenjunga and Mt. Everest\,:} Estimate the relative intensity of Mt
Kanchenjunga and Mt Everest with the above level of pollution with respect to
the reference value. Which, if any of these peaks will be visible from
Darjeeling? Assume that pollution is uniform throughout the path the light
travels. \hfill 1pt
